\originalchapter{从原子状态到恒星大气的光}
本章支撑8.901第4--6次课。第\ref{ch:radiation}章已给比强度、Planck函数、输运方程和Saha推导；这里把这些关系接到实际吸收线、组成、不透明度和表面边界，说明何时可以把光谱当作温度计。

\originalsection{电离与激发分别决定一条谱线能否出现}
同一离子内部的激发遵循Boltzmann比$n_j/n_i=(g_j/g_i)\exp[-(E_j-E_i)/(k_BT)]$；不同电离级的比例遵循Saha式\eqref{eq:saha}，还取决于自由电子密度。某条吸收线需要低能级有人占据，也需要对应物种尚未全部电离。因此线强常在中间温区达到最大，不能认为温度越高所有高能线都越强。

以氢Balmer线为例，下能级$n=2$在低温很少，升温可增加它；更高温时中性氢减少，又削弱线。恒星谱型序列主要反映有效温度，密度、重力、金属丰度和辐照也改变线。Saha只适合局部热平衡、非简并、碰撞足够的气体；星际电离区可能主要由紫外光照建立电离平衡，此时应比较光电离和复合率，不能只代入一个低密度温度。

多个电离级须同时满足核素总数与电荷守恒。例如氦$x_0+x_1+x_2=1$、电子份额$x_e=x_1+2x_2$，两条Saha关系中的$n_e$由这组未知量共同决定。PS2.3的实际求根提供完整例子；把$n_e$事先固定为“全电离值”会在部分电离温区破坏守恒。

\originalsection{平均粒子质量与压强来源}
理想气体压强$P_{\rm gas}=\rho k_BT/(\mu m_u)$，$\mu$是每个可自由平动粒子的平均重子质量数。完全电离的混合气体，质量分数$X_i$、原子质量数$A_i$、电荷$Z_i$给
\[
 \mu^{-1}=\sum_i\frac{X_i(1+Z_i)}{A_i},\qquad
 \mu_e^{-1}=\sum_i\frac{X_iZ_i}{A_i}.
\]
前式同时数离子、电子，后式只数电子。纯氢完全电离$\mu=1/2,\mu_e=1$；纯氦$\mu=4/3,\mu_e=2$。中性和部分电离时要按实际粒子数重算。使用$m_p$还是$m_u$为归一也应一致，精密核质量计算不能把它们完全等同。

辐射压强$P_{\rm rad}=a_RT^4/3$，简并电子压强则由填充的动量态产生。恒星总压强可以同时有这些项。对流和稳定性所用$\Gamma_1=(\partial\ln P/\partial\ln\rho)_s$不是总能处处等于理想单原子气体的$5/3$；电离可吸收压缩热，辐射或相对论电子也使有效响应变软。

\originalsection{不透明度是每单位质量的消光能力}
定义$\alpha_\nu=\rho\kappa_\nu$，平均自由程$\ell_\nu=1/\alpha_\nu$。$\kappa$单位面积／质量，不等于光学深度。沿射线$\tau_\nu=\int\rho\kappa_\nu\dd s$，同样的$\kappa$在不同密度、几何长度下可给完全不同的透明度。

电子散射、束缚--束缚谱线、束缚--自由电离、自由--自由辐射及低温气体的负氢离子是不同机制。完全电离时Thomson散射近似频率无关，按成分$\kappa_{\rm es}=n_e\sigma_T/\rho$。常用Kramers连续吸收$\kappa\propto\rho T^{-7/2}$是有限温区、组成和频段平均下的近似；强谱线、部分电离、金属和分子可改变它，不能全星强行用一个幂律。

光学厚扩散时第\ref{ch:radiation}章的Rosseland输运推导使透明窗口权重更大；光学薄辐射损失则偏向发射强的频段，使用不同平均。散射会改变方向而未必热化光子，若真吸收不足，不能由“总光学深度大”直接推出局部黑体。

\originalsection{灰大气给出一个可追溯的表面条件}
取平面、稳态、灰、LTE大气，光学深度$\tau$向内增加，比强度对外向余弦$\mu_\theta$。定义辐射矩
\[
 J=\frac12\int_{-1}^1I\dd\mu_\theta,\quad
 H=\frac12\int_{-1}^1\mu_\theta I\dd\mu_\theta,\quad
 K=\frac12\int_{-1}^1\mu_\theta^2I\dd\mu_\theta .
\]
其中$F=4\pi H$。输运方程$\mu_\theta\dd I/\dd\tau=I-S$积分给$\dd H/\dd\tau=J-S$、$\dd K/\dd\tau=H$。辐射平衡要求$J=S=B$，所以$H$常数；采用Eddington闭合$K=J/3$得$J=3H\tau+C$。

再用近似无外照边界$J(0)=2H$，以及积分黑体$B=\sigma_{\rm SB}T^4/\pi$、$F=\sigma_{\rm SB}T_{\rm eff}^4$，得到
\[
 T^4(\tau)=\frac34T_{\rm eff}^4(\tau+2/3).
\]
因此$T(\tau=2/3)=T_{\rm eff}$；这个$2/3$来自近似边界和矩闭合，不是任何真实恒星都在同一个频率独立层发光。斜射线看到较浅、较冷层，产生边缘昏暗的方向：形式解$I(0,\mu_\theta)=\int_0^\infty S(\tau)e^{-\tau/\mu_\theta}\dd\tau/\mu_\theta$明确展示斜线的权重。

大气与内部扩散在适当深度匹配才完成恒星边界问题。直接令表面$T=0$虽方便求多方球的数学表面，却不能同时作为真实发射光谱的边界条件。
